\documentclass[11pt]{article}
\usepackage[T1]{fontenc}
\usepackage{lmodern}
\usepackage[margin=1in]{geometry}
\usepackage{amsmath,amssymb,amsthm,microtype}
\usepackage[hidelinks]{hyperref}

\newtheorem{theorem}{Theorem}
\title{A gap above the \(n^2/3\) bound for maximal orthogonal partial Latin squares}
\author{Aafi Malek}
\date{26 November 2026}

\begin{document}
\maketitle

\begin{abstract}
Let \(F\) be the number of filled cells in a maximal pair of orthogonal
partial Latin squares of order \(n=3s\). The lower bound \(F\ge 3s^2\)
is known. We exclude \(F=3s^2+1\) for \(s\ge2\). For \(s\ge3\), we
show that \(F=3s^2+2\) occurs precisely in unions of three full
orthogonal pairs of orders \(s-1,s,s+1\). At order \(18\) this gives
a lower bound of \(111\) filled cells.
\end{abstract}

\section{Statement and notation}

Represent two superimposed partial Latin squares by a set \(D\) of
words of length four. The four coordinates describe the row, column
and two symbols. Each coordinate has its own set of \(n\) labels.
The Latin and orthogonality conditions say that any two distinct words
agree in at most one coordinate. Call a pair of labels in distinct
coordinates \emph{used} if it occurs in some word. The pair is maximal
if every possible word contains a used pair. The number of filled
cells is \(F=|D|\).

Write \(f(x)\) for the frequency of a label \(x\). In each other
coordinate, \(x\) has exactly \(f(x)\) used neighbors, because no pair
can occur in two words. The frequencies sum to \(F\) in each coordinate.
A full orthogonal pair of order \(b\) uses every pair between any two
of its four \(b\)-element label sets.

\begin{theorem}\label{thm:main}
Let \(D\) be maximal of order \(n=3s\), where \(s\ge2\).
Then \(F=3s^2+1\) is impossible. If \(s\ge3\), then \(F=3s^2+2\)
holds precisely when \(D\) consists of three full orthogonal pairs,
on disjoint label sets in every coordinate, of orders \(s-1,s,s+1\).
\end{theorem}

\section{The frequency bound}

We recall the matching bound of Donovan, Grannell and
Yaz\i c\i~\cite[Theorem 2.1 and Lemma 2.2]{feb}. The equality conditions
will be needed. Put \(m=\min_x f(x)\). Fix any label of frequency \(m\)
and write \(q=n-m\). In each of the other coordinates there are exactly
\(q\) labels unused with the fixed label. Choose any two of those
coordinates and let \(t\) be the size of a maximum matching of unused
pairs between the two sets of \(q\) labels. For each matched pair
\((x,y)\), maximality requires that every one of the \(q\) candidate
labels in the remaining coordinate is used with \(x\) or \(y\).
Consequently \(f(x)+f(y)\ge q\).

For any two unmatched labels on opposite sides, their pair is used,
and all pairs between the two unmatched sets are used. For each
matching edge, at least one of the two pairs joining its endpoints to
the two unmatched labels is used. Otherwise replacing the matching
edge by those two unused pairs would increase the matching size.
The two unmatched labels therefore have frequency sum at least
\(2(q-t)+t=2q-t\).

Pair the unmatched labels arbitrarily between the two coordinates.
Also pair the \(m\) excluded labels on each side, each of which has
frequency at least \(m\). Summing the three types of frequency sums
gives
\begin{equation}\label{eq:bound}
2F\ge 2m^2+tq+(q-t)(2q-t)
=2m^2+q^2+(q-t)^2
=6s^2+3(s-m)^2+(q-t)^2.
\end{equation}
If equality holds, each of the \(m\) excluded labels on either side
has frequency \(m\). The bound and this consequence hold for any
choice of the fixed minimum-frequency label and of the two
coordinates.

\section[Minimum frequency s]{The case \(m=s\)}

Put \(k=F-3s^2\), and suppose \(0\le k\le2\) and \(k<s\).
Call frequency-\(s\) labels low and the others high. Each coordinate
has at most \(k\) high labels because its total excess over frequency
\(s\) is \(k\).

\paragraph{Common neighbors.}
Two low labels in distinct coordinates that share a used neighbor
must form a used pair. To see this, suppose their pair is unused.
In each of the remaining two coordinates, the set of labels unused
with both low labels has size at least \(n-2s=s\). If the low labels
share a neighbor, one of the two candidate sets has size at least
\(s+1\). By maximality all pairs between the candidate sets are used.
This forces every label of the other candidate set to be high, giving
at least \(s\) high labels in one coordinate, contrary to \(k<s\).

\paragraph{Components of low labels.}
Take the graph on low labels whose edges are the used pairs. Every
vertex has a low neighbor in each other coordinate, since it has
\(s\) used neighbors there and at most \(k<s\) are high. If two low
labels \(x,x'\) in the same coordinate share a low neighbor \(y\),
they have identical low neighborhoods. Indeed, choose a low neighbor
\(w\) of \(x\) in a third coordinate. The common-neighbor property
gives \(yw\) through \(x\) and \(x'w\) through \(y\). For a low neighbor
\(v\) of \(x\) outside \(y\)'s coordinate, the same property gives
\(vy\) through \(x\) and \(vx'\) through \(y\). For a low neighbor
\(v\) in \(y\)'s coordinate, use \(w\) in place of \(y\).
Interchanging \(x,x'\) gives the reverse inclusion.

Every component is complete across its coordinate parts. In a
shortest path of length at least three, the first and third vertices
either lie in distinct coordinates and are adjacent by the
common-neighbor property, or lie in the same coordinate and have
identical low neighborhoods. In the latter case the first vertex
is adjacent to the fourth. Both cases shorten the path. A path of
length two between distinct coordinates can also be shortened.
Every component meets all four coordinates.

\paragraph{Adding the high labels.}
Each high label has a low neighbor in every other coordinate.
All its low neighbors lie in one component, since those in distinct
coordinates are adjacent through the high label, and those in the
same coordinate are joined through a low neighbor of the high label
in another coordinate. Assign
the high label to that component.

A low vertex has all of its \(s\) used neighbors in each other
coordinate inside its assigned component. Each component therefore
has at least \(s\) labels in each coordinate after assignment.
Before assignment, a component has at most \(s\) low labels in
each coordinate because all of them neighbor any low vertex of
that component in another coordinate. Each coordinate has at least
\(3s-k>2s\) low labels. There are thus at least three components.
The lower bound of \(s\) labels per component and coordinate permits
at most three. It follows that there are exactly three components
with exactly \(s\) labels in each coordinate.

\paragraph{Words across components.}
Any word that meets more than one component consists entirely of
high labels. If \(k>0\), there is at least one such word, since at most
\(3s^2\) words lie entirely within components. Any word meeting more
than one component still has two coordinates from the same component,
because it has four coordinates and there are three components.
Its pair in those coordinates cannot occur in an internal word.
Counting that projection gives at most \(3s^2-1\) internal words.
Thus there are at least \(k+1\) words meeting more than one component.

For \(k=1\), there is at most one all-high word. For \(k=2\), there
are at most two all-high words. Among three such words, each coordinate
would have an agreement between at least two of them, since it uses
at most two high labels. Four coordinates and three pairs of words
would force some two words to agree twice. These upper bounds
contradict the required \(k+1\). For \(k=0\) every word lies within
a component, and counting pairs shows that all three components are
full orthogonal pairs of order \(s\).

\section[Minimum frequency s minus 1]{The case \(m=s-1\) at \(F=3s^2+2\)}

For \(F=3s^2+1\), the average frequency and \eqref{eq:bound}
force \(m=s\), so the previous section proves the first assertion
of Theorem~\ref{thm:main}. When \(F=3s^2+2\) and \(s\ge3\), the
same bounds give \(m\in\{s-1,s\}\). The case \(m=s\) was just excluded.

Suppose \(m=s-1\), so that \(q=2s+1\). Substitution in
\eqref{eq:bound} gives \((q-t)^2\le1\). If \(t=q\), then each of
the \(q\) candidate labels in the fourth coordinate must be used
with at least one endpoint of each matching edge. A label of
frequency \(f\) can meet at most \(2f\) such edges. All \(q\)
candidates would have frequency at least \(s+1\), giving
\[
F\ge (s-1)^2+(2s+1)(s+1)=3s^2+s+2,
\]
a contradiction. Therefore \(t=q-1\), and \eqref{eq:bound} is
tight. Its equality condition says that every used neighbor of a
minimum-frequency label is also of minimum frequency, since the
two coordinates in the matching argument can be chosen freely.

Any two minimum-frequency labels from distinct coordinates form a
used pair. Otherwise in each remaining coordinate at least
\(3s-2(s-1)=s+2\) labels are unused with both. All pairs between
these two candidate sets must be used. In one coordinate this gives
at least \(s+2\) labels of frequency at least \(s+2\), so
\(F\ge 3s(s-1)+3(s+2)=3s^2+6\), again a contradiction.

There is at least one minimum-frequency label in each coordinate,
because \(m\ge1\) and all neighbors of a minimum-frequency label
also have frequency \(m\). The preceding two properties show that
each coordinate has exactly \(m\) such labels. All words meeting
these labels form a full orthogonal pair of order \(m=s-1\), with
no used pair from those labels to the remaining ones.

The remaining words have \(q=2s+1\) labels per coordinate and
\(F-m^2=s^2+(s+1)^2=\lceil q^2/2\rceil\) words. Their projection onto
any three coordinates is a maximal partial Latin square. Indeed,
if three remaining labels had all their pairs unused, adjoining
a minimum-frequency label in the fourth coordinate would give a
word with every pair unused in \(D\). Also, no two remaining words
have the same projection, by pair injectivity.

By Donovan, Grannell and Yaz\i c\i~\cite[Lemma 1.1]{feb}, each
such projected square splits into full Latin squares of orders \(s\)
and \(s+1\). The two blocks are distinguished by label frequency in
the remaining words, which is preserved by projection. Using the
projections onto coordinates \(1,2,3\) and \(1,2,4\), every remaining
word consists entirely of labels of one of those two frequencies.
There are \(s^2\) words in the first class and \((s+1)^2\) in the
second. Pair injectivity makes each class a full orthogonal pair.
This proves the classification.

Conversely, disjoint full pairs of orders \(s-1,s,s+1\) have
\(3s^2+2\) words. The union is maximal because every four-coordinate
word has at least two labels from the same block, whose pair is
used there. \qed

\section{Consequences and prior work}

Full orthogonal pairs exist at every positive order except \(2\)
and \(6\) \cite[Section 1]{feb}. In particular, the second-smallest
attainable size at order \(3s\) is \(3s^2+2\) for \(s=4\) and every
\(s\ge8\). The order \(12\) example uses blocks of orders \(3,4,5\)
and has \(50\) filled cells, two above the minimum of \(48\).

At order \(18\), the case \(F=108\) would require an orthogonal
pair of order \(6\), by \eqref{eq:bound} and the \(m=s,\ k=0\)
argument. The theorem excludes \(109\) and \(110\), the latter
because the classified orders would be \(5,6,7\). Full blocks of
orders \(4,7,7\) give \(114\) filled cells, so
\(111\le \min F\le114\). At order \(6\), the analogous equality
case needs order \(2\) blocks, and the theorem excludes \(13\).
Together with the \(3,3\) construction this gives
\(14\le \min F\le18\). These orders were among the unresolved
values in Donovan, Grannell and Yaz\i c\i~\cite[Table 1]{original}.

The February and November 2026 preprints give the universal bound
and classify the arrangements attaining it at large orders
\cite{feb,nov}. Horak and Kim~\cite{room} mention work in preparation
by Hor\'ak, Meszka and Rosa on the spectrum of attainable sizes.
We found no public proof of the two claims in Theorem~\ref{thm:main}
in the sources searched. Priority remains open to correction.
The accompanying script checks the finite constructions by testing
all pair projections and all potential insertions. The proof above
establishes the general exclusions and classification.

\begin{thebibliography}{9}
\bibitem{original} D. M. Donovan, M. J. Grannell and E. \c{S}. Yaz\i c\i,
\emph{On maximal orthogonal partial Latin squares and minimal codes with
specified length, minimum distance and covering radius},
Des. Codes Cryptogr. \textbf{93} (2025), 5099--5113.
\href{https://doi.org/10.1007/s10623-025-01704-x}{doi:10.1007/s10623-025-01704-x}.

\bibitem{feb} D. M. Donovan, M. Grannell and E. \c{S}. Yaz\i c\i,
\emph{On the minimum number of entries in a pair of maximal orthogonal
partial Latin squares}, \href{https://arxiv.org/abs/2602.09908}{arXiv:2602.09908} (2026).

\bibitem{nov} D. M. Donovan, M. J. Grannell and E. \c{S}. Yaz\i c\i,
\emph{On the minimum number of entries in a set of maximal mutually
orthogonal partial Latin squares},
\href{https://arxiv.org/abs/2611.13961}{arXiv:2611.13961} (2026).

\bibitem{room} P. Horak and D. Kim,
\emph{An improved lower bound for maximal partial Room squares},
\href{https://arxiv.org/abs/2611.07117}{arXiv:2611.07117} (2026).
\end{thebibliography}

\noindent\small AI assistance was used to develop the proofs, search the
literature, write the note, and implement the witness checker.
\end{document}
